\section{\ref{sec:dendrites}장. 가지돌기 개수}

부류의 구조와 입력 개수는 해부학과 전기 기록으로 측정되었다. 추정~\eqref{eq:dendriteload}은 커패시터의 간단한 물리이고, 입력 개수에 대한 계산적 설명은 이론과 모델에 기댄다.

{\footnotesize
\setlength{\tabcolsep}{3pt}\renewcommand{\arraystretch}{1.15}
\begin{longtable}{>{\raggedright}p{0.5cm}>{\raggedright}p{4.0cm}>{\raggedright}p{5.6cm}>{\raggedright}p{2.3cm}>{\raggedright\arraybackslash}p{1.7cm}}
\toprule
No. & 명제 & 실험과 측정한 것 & 출처 & 상태 \\
\midrule\endfirsthead
\toprule
No. & 명제 & 실험과 측정한 것 & 출처 & 상태 \\
\midrule\endhead
1 & 막의 단위 면적당 커패시턴스 $c_m\approx1$~\textmu F/cm$^2$ & 뉴런 세포체에서 막을 피펫으로 떼어 내고 전압 계단을 걸어, 용량성 전류의 감쇠로 전체 커패시턴스를 구한 뒤 면적으로 나눴다: 세 부류의 뉴런에서 $0.9$~\textmu F/cm$^2$ & \cite{Gentet2000} & 측정됨 \\
2 & 충전 시간 $t\approx c_mA\Delta V/I$~\eqref{eq:dendriteload}: 큰 뉴런에는 동시 입력 수십 개가 필요하고, 작은 뉴런에는 큰 시냅스 하나면 충분하다 & 커패시터 충전 법칙에 의한 추정. 수치($20$과 $300$~pF, $0.1$과 몇 nA)는 크기 수준 & 본문에서 유도 & 이론 \\
3 & 거대 시냅스 하나는 몇 나노암페어의 전류를 주어 작은 뉴런을 곧바로 문턱값까지 올린다 & 절편에서 말단과 수신 세포를 동시 기록: 시냅스 하나의 전류 $2$--$13$~nA, 입력 스파이크마다 문턱 이상의 반응, $36^\circ$C에서 지연 약 $0.4$~ms; 수신 세포는 \nt{GLYC}를 낸다 & \cite{Borst1995,BorstSoria2012} & 측정됨 \\
4 & 가지돌기 0개: 촉각과 통증의 감각 뉴런에서 스파이크는 세포체를 거치지 않고 센서에서 척수로 간다 & 구조(세포체 옆에서 둘로 갈라지는 축삭)는 해부학으로 확립되었다. 전도에 세포체의 흥분성이 필요 없다는 것은 이런 뉴런의 검증된 모델에서 보였다: 흥분성 세포체가 없어도 스파이크는 분기점을 똑같이 확실하게 통과한다 & \cite{AmirDevor2003} & 이론 \\
5 & 가지돌기 1개: 후각 뉴런은 수용체 한 종류를 지니고 입력 선로 하나를 전달한다 & 수용체 유전자 실험의 리뷰: 후각 뉴런 하나는 큰 유전자군 가운데 수용체 유전자 하나를 발현한다 & \cite{Mombaerts2004} & 측정됨 \\
6 & 가지돌기 2개: 소리 방향 검출기에서 왼쪽 귀와 오른쪽 귀의 입력은 서로 다른 가지돌기에 붙는다 & 리뷰에 모은 포유류의 해부학과 기록: 두 귀의 입력이 반대쪽을 향한 두 가지돌기에 나뉘어 있다 & \cite{Grothe2010} & 측정됨 \\
7 & 입력을 두 가지돌기에 나누면 AND가 더 엄격해진다 & 모델에서 입증: 한 가지돌기의 포화~\eqref{eq:shunt} 때문에 한쪽 입력만으로는 문턱값에 이르지 못하고, 동시성 검출이 개선된다. 입력 배치를 바꾼 직접 실험은 없다 & \cite{AgmonSnir1998} & 이론 \\
8 & 운동 학습 회로에 작은 \nt{GLUT} 뉴런 약 $7\cdot10^{10}$개, 각각 입력 $\sim4$개 & 등방성 분획법에 의한 세포 계수: 사람 뇌의 이 부분에 뉴런 약 $690$억 개. 살아 있는 쥐에서 이런 세포 하나를 기록: 촉각에 소수의 입력에서 오는 이산 전류의 버스트가 도착하고, 이들이 더해지면 세포가 발화한다 & \cite{Azevedo2009,Chadderton2004} & 측정됨 \\
9 & 입력 $n$개인 문턱 소자는 무작위 패턴을 최대 $P_{\max}\approx2n$개 분리한다~\eqref{eq:cover}; 운동 학습 회로의 출력 뉴런에서 이 경계는 $\sim2\cdot10^5$, 현실적인 추정은 $\sim4\cdot10^4$개의 대응 & 경계는 조합 기하학의 고전적 결과. 현실적인 추정은 양의 가중치만 쓰고 잡음 여유를 둔 용량 이론을, 이 뉴런의 측정된 입력 수에 적용한 것이다 & \cite{Cover1965,Brunel2004} & 이론 \\
10 & 입력이 적은 믹서는 입력을 확장하려고 필요하다. “네 개”는 최적에 가깝다 & 이론: 믹서 층이 주는 표현의 차원은 해부학적으로 관찰되는 입력 개수 근처에서 최대가 된다. 확장이라는 원래 가설은 고전 연구에 있다 & \cite{Marr1969,Albus1971,LitwinKumar2017} & 가설 \\
11 & 큰 나무: 입력 $10^4$--$10^5$개 & 전자 현미경 사진으로 시냅스 계수: 쥐의 운동 학습 회로 출력 \nt{GABA} 뉴런 하나에 시냅스 $\approx1.7\cdot10^5$개; 해마 \nt{GLUT} 뉴런에 흥분성 입력 약 $30\,000$개와 억제성 입력 $1\,700$개 & \cite{NapperHarvey1988,Megias2001} & 측정됨 \\
12 & 큰 나무의 가지는 각자 문턱값을 가진 별도의 하위 회로이며, 뉴런은 작은 다층 네트워크다 & 가는 가지돌기의 두 지점을 국소 자극: 같은 가지의 입력은 시그모이드로, 다른 가지의 입력은 선형으로 더해진다. 사람 피질 뉴런의 가지돌기에서 세포 하나가 선형 분리 불가능한 문제를 풀게 하는 등급형 칼슘 스파이크가 발견되었다. 모델: 뉴런 하나를 재현하려면 깊은 네트워크가 필요하다 & \cite{Polsky2004,Gidon2020,Beniaguev2021} & 측정됨 \\
13 & 출력 가지돌기: 망막 \nt{GABA} 뉴런의 가지 하나하나가 별도의 방향 검출기다 & 개별 가지의 2광자 칼슘 기록: 가지의 신호는 세포체에서 끝 쪽으로 가는 움직임에 선택적이지만 세포체의 전압은 그렇지 않다 & \cite{Euler2002} & 측정됨 \\
14 & 입력 개수는 과제를 따른다(명제~\ref{prop:dendrites}) & 부류의 구조는 측정되었다(3--13행). 입력 개수가 속도나 분류를 위해 선택되었다는 것은 개별 부류에 대해 이론과 모델로만 보였다 & \cite{LitwinKumar2017,AgmonSnir1998} & 가설 \\
\bottomrule
\end{longtable}}
